\paragraph{C.3x (Proved/refuted): the complete ordinary rational degree-one dimension table in positive characteristic.}
Fix an algebraically closed field $k$ of characteristic $p>0$.
Let $R$ be the strict henselization at a $k$-rational point of a
normal local scheme essentially of finite type over $k$, with the
given coefficient field inducing the residue-field identification.
Write
\[
 N=\dim R,\qquad
 X=\operatorname{Spec}R\setminus\{\mathfrak m\},\qquad
 Y=X\cap V(f),\quad 0\ne f\in\mathfrak m,
\]
and assume $X$ and $Y$ regular. Use ordinary algebraic equivalence,
defined by actual line-bundle families on products with smooth
connected curves, and put
\[
 J^1(X)=\operatorname{Pic}(X)_{\mathrm{alg}}\otimes_{\mathbb Z}\mathbb Q,
 \qquad
 J^1(Y)=\operatorname{Pic}(Y)_{\mathrm{alg}}\otimes_{\mathbb Z}\mathbb Q.
\]
If $k=\overline{\mathbb F}_p$, both groups vanish in every dimension.
If $k\ne\overline{\mathbb F}_p$, the universal assertions in this
algebraizable scope have exactly the following answers.
Here ``No'' means that the explicit example below over the fixed
field $k$ disproves the universal assertion.

\begin{center}
\begin{tabular}{lll}
\toprule
Dimension $N$ & Always injective & Always surjective \\
\midrule
$1$ & Yes, both groups vanish & Yes, both groups vanish \\
$2$ & No, the elliptic cone below & Yes, the target vanishes \\
$3$ & Yes, by C.3n & No, the same cone as divisor below \\
$\ge4$ & Yes, by C.3n & Yes, by C.3w \\
\bottomrule
\end{tabular}
\end{center}

In particular, the two boundary failures hold in characteristic
$2$, in characteristic $3$, and in every other positive
characteristic. The last row uses the algebraizability hypothesis
of C.3w. It is not an assertion for arbitrary nonalgebraizable
excellent germs, and it differs from the characteristic-zero
surjectivity table in C.3o.

\begin{proof}
\emph{1. A smooth constant elliptic curve in every characteristic.}
Use homogeneous coordinates $(x:y:z)$ and define
\[
 F_p=
 \begin{cases}
   y^2z+yz^2+x^3,&p=2,\\
   y^2z-x^3+xz^2,&p\ne2,
 \end{cases}
 \qquad
 E=V_+(F_p)\subset\mathbb P^2_k,
 \qquad O=(0:1:0).
\]
For $p=2$, the three partial derivatives are $x^2,z^2,y^2$.
They have no common projective zero. For $p\ne2$, the derivatives
are
\[
 (F_x,F_y,F_z)=(z^2-3x^2,\ 2yz,\ y^2+2xz).
\]
At $z=0$ the equation forces $x=0$, and then $F_z=y^2\ne0$.
On $z=1$, a singular point would have $y=0$ and
$x(x^2-1)=0$. The value of $F_x$ is $1$ when $x=0$ and $-2$
when $x=\pm1$, so it cannot vanish. This also checks $p=3$.
An exact identity over $\mathbb Z[1/2][x,y]$, with
$G=y^2-x^3+x$, certifies the same affine assertion:
\begin{equation}\label{eq:C3x-odd-smoothness}
 1=
 \left(1-\frac{3x^2}{2}\right)G_x
 +\frac{9x}{2}G
 -\frac{9xy}{4}G_y.
\end{equation}
Thus $E$ is a smooth plane cubic with the specified point $O$,
and hence a genus-one curve with its elliptic-curve structure.
All these data are defined over $\mathbb F_p$.

\emph{2. A specified infinite-order actual class over the fixed field.}
Let $k_0\subset k$ be the algebraic closure of $\mathbb F_p$ in $k$.
The assumption $k\ne k_0$ gives an element $t\in k$
transcendental over $k_0$. Choose $v\in k$ satisfying
\[
 v^2+v=t^3\quad(p=2),\qquad
 v^2=t^3-t\quad(p\ne2),
\]
and set
\[
 P=(t:v:1)\in E(k),\qquad Q=\mathcal O_E(P-O).
\]
The Abel map $E\to\operatorname{Pic}^0(E)$, $P\mapsto
\mathcal O_E(P-O)$, is an isomorphism over $k_0$.
For example, Riemann--Roch shows that each degree-one line
bundle on a genus-one curve has a unique effective divisor;
this gives the inverse. For every integer $a>0$, the group
scheme $E[a]$ is finite over $k_0$, including when $p\mid a$:
multiplication by $a$ is a finite flat isogeny of degree $a^2$
\citep[Proposition 8.1 and Theorem 8.2, pp.~12--14]{D_MilneAbelian2022}.
Every $k$-point of a finite $k_0$-scheme is already a $k_0$-point,
since its residue field is finite over the algebraically closed
field $k_0$. Consequently all torsion points of this constant
elliptic curve over $k$ lie in $E(k_0)$. The point $P$ does not,
because its $x$-coordinate is $t$. Hence $Q$ has infinite order.
This argument includes the non-etale $p$-primary torsion; an
etale-only torsion calculation would not suffice.

\emph{3. The actual Picard injection on the strict local cone.}
Put $L=\mathcal O_E(1)$, so $\deg L=3$, and define
\[
 R_2=\left(k[x,y,z]/(F_p)\right)_{(x,y,z)}^{\mathrm{sh}},
 \qquad
 X_2=\operatorname{Spec}R_2\setminus\{\mathfrak m_2\}.
\]
The cone is integral and smooth away from its vertex.
It is normal: the hypersurface is Cohen--Macaulay and is
regular in codimension one. Its local ring is excellent,
and its strict henselization remains an excellent normal
local domain. The usual cone projection induces
$\rho:X_2\to E$. Base change the proper cone resolution
$\operatorname{Tot}(L^{-1})\to\operatorname{Spec}k[x,y,z]/(F_p)$
to $\operatorname{Spec}R_2$. The resulting regular integral
scheme $Z_2$ is isomorphic to $X_2$ away from its single
exceptional prime divisor $E$, and
\[
                     \mathcal O_{Z_2}(E)|_E=L^{-1}.
\]
If $Q'\in\operatorname{Pic}^0(E)(k)$ pulls back to a trivial
line bundle on $X_2$, its pullback to $Z_2$ has a rational
trivialization whose Cartier divisor is supported on $E$.
It is therefore $\mathcal O_{Z_2}(bE)$ for some integer $b$.
Restriction gives $Q'\simeq L^{-b}$, and degrees give
$0=-3b$, so $b=0$ and $Q'=0$. This proves an actual injection
\begin{equation}\label{eq:C3x-elliptic-injection}
 \rho^*:\operatorname{Pic}^0(E)(k)
       \lhook\joinrel\longrightarrow
       \operatorname{Pic}(X_2)_{\mathrm{alg}}.
\end{equation}
The target subscript is justified directly, not by an
identity-component convention for a local Picard functor:
the line bundle $\mathcal O_{E\times E}
(\Delta-\{O\}\times E)$ has fibres
$\mathcal O_E(P'-O)$ as $P'\in E(k)$ varies. Pull it back
along $\rho\times\mathrm{id}_E$ to obtain an actual family
on $X_2\times E$. In particular, $\rho^*Q$ has infinite
order and gives a nonzero element of $J^1(X_2)$.

\emph{4. Failure of injection in dimension two.}
Take the specified nonzero divisor
\[
 s=
 \begin{cases}x,&p=2,\\y,&p\ne2\end{cases}
 \quad\text{in }R_2,
 \qquad Y_2=X_2\cap V(s).
\]
Its equation on the plane section factors as
\[
 F_2|_{x=0}=yz(y+z),\qquad
 F_p|_{y=0}=-x(x-z)(x+z)\quad(p\ne2).
\]
In each case these are three distinct lines meeting only
at the vertex. Removing that vertex gives a regular
zero-dimensional puncture, and strict henselization
preserves this regularity. More explicitly, normalization
has three strictly henselian discrete valuation local
factors, so $Y_2$ is a disjoint union of three spectra
of fields. Thus $\operatorname{Pic}(Y_2)=0$.
The nonzero rational class $\rho^*Q\otimes1$ lies in the
kernel of $J^1(X_2)\to J^1(Y_2)$.

\emph{5. Failure of surjection in dimension three.}
Define instead
\[
 R_3=k[x,y,z]_{(x,y,z)}^{\mathrm{sh}},\qquad
 X_3=\operatorname{Spec}R_3\setminus\{\mathfrak m_3\},
 \qquad f=F_p.
\]
The ring $R_3$ is regular local, and its puncture has
trivial Picard group
\citep[Lemma 52.27.5, Tag 0F2H]{D_StacksRegularPuncture}.
Strict henselization commutes with the quotient by
$F_p$, so $Y_3=X_3\cap V(F_p)$ is the preceding $X_2$.
It is regular, and the nonzero element
$\rho^*Q\otimes1\in J^1(Y_3)$ has no preimage from
$J^1(X_3)=0$.

\emph{6. Completion of the universal table and the finite-field case.}
For $N=1$, a normal noetherian local domain is a discrete
valuation ring; $X$ is a field spectrum and $Y$ is empty,
so both groups vanish. For $N=2$, regularity makes $Y$
a zero-dimensional noetherian scheme, hence a finite
disjoint union of spectra of fields. Its Picard group
vanishes, proving universal surjectivity in that row.
For $N\ge3$, C.3n proves that ordinary Picard restriction
has only $p$-primary torsion in its kernel, and therefore
is injective after tensoring with $\mathbb Q$. For
$N\ge4$, C.3w proves integral ordinary surjectivity in
exactly the algebraizable scope stated here. Tensoring
with $\mathbb Q$ gives the asserted isomorphism.

When $k=\overline{\mathbb F}_p$, every point of every
abelian variety over $k$ is torsion: the variety, its
group law, and the given point descend to some finite
field, whose point group is finite. The actual
dual-Jacobian-family construction of C.3n places
each ordinarily algebraically trivial line-bundle
class in an abelian point-group image. Thus every
such class is torsion and both rational groups vanish.
This last conclusion does not require algebraizability;
it has the regular noetherian coefficient-field scope
of that actual-family argument.
\end{proof}

The exact certificate in
\texttt{code/verify\_positive\_characteristic\_boundary.py}
checks the polynomial identity
\eqref{eq:C3x-odd-smoothness} over $\mathbb Q$, its sole
denominator prime $2$, the characteristic-two derivative
identities, and the two squarefree hyperplane factorizations.
It also independently enumerates the projective singular
and section points over the prime fields of characteristic
$2,3,5,7,11,13$. The universal smoothness and infinite-order
assertions are proved above; they are not inferred from
this finite sample.
